\begin{center}
\begin{tikzpicture}[x=.116\linewidth,y=7mm,every node/.style={font=\scriptsize,align=center}]
\fill[chFourTint] (0,3) rectangle (8,5);
\draw (0,0) rectangle (8,5);
\draw (2.2,0)--(2.2,5) (3.6,0)--(3.6,4) (5,0)--(5,4) (6.2,0)--(6.2,5);
\draw (0,3)--(8,3) (2.2,4)--(6.2,4) (6.2,1)--(8,1) (6.2,2)--(8,2);
\node at (1.1,4) {Student\\Name};\node at (4.2,4.5) {Compulsory};\node at (7.1,4) {Optional};
\node at (2.9,3.5) {Bangla};\node at (4.3,3.5) {English};\node at (5.6,3.5) {ICT};
\node at (1.1,1.5) {Mahbub};
\node at (2.9,1.5) {$\checkmark$};\node at (4.3,1.5) {$\checkmark$};\node at (5.6,1.5) {$\checkmark$};
\node at (7.1,2.5) {Physics};\node at (7.1,1.5) {Math};\node at (7.1,.5) {Biology};
\end{tikzpicture}
\captionof{figure}{অনুশীলন 33: শিরোনাম ও তথ্যের সারিতে ঘরের বিস্তার}
\end{center}
\textbf{চিত্র ব্যাখ্যা:} Student Name ও Optional শিরোনাম দুই সারি জুড়ে আছে। Compulsory শিরোনাম তিন কলাম জুড়ে আছে। তথ্যের অংশে Mahbub ও তিনটি টিকচিহ্নের ঘর তিন সারি জুড়ে; পাশে Physics, Math ও Biology আলাদা তথ্যের সারিতে আছে।
